\documentclass[11pt]{article}
\usepackage[letterpaper,margin=0.83in,headheight=16pt]{geometry}
\usepackage[T1]{fontenc}
\usepackage[utf8]{inputenc}
\usepackage{lmodern}
\usepackage{microtype}
\usepackage{amsmath,amssymb}
\usepackage{xcolor}
\usepackage{hyperref}
\usepackage{enumitem}
\usepackage{booktabs,longtable,tabularx,array}
\usepackage{fancyhdr}
\usepackage{titlesec}
\usepackage{listings}
\usepackage{tikz}
\usepackage{float}
\definecolor{suiteblue}{HTML}{17365D}
\definecolor{suitelight}{HTML}{EAF1F8}
\definecolor{suitegray}{HTML}{556070}
\hypersetup{colorlinks=true,linkcolor=suiteblue,urlcolor=suiteblue,breaklinks=true}
\titleformat{\section}{\Large\bfseries\color{suiteblue}}{\thesection}{0.55em}{}
\titleformat{\subsection}{\large\bfseries\color{suiteblue}}{\thesubsection}{0.55em}{}
\setlist[itemize]{itemsep=0.25em,topsep=0.35em,leftmargin=1.4em}
\setlist[enumerate]{itemsep=0.25em,topsep=0.35em,leftmargin=1.6em}
\setlength{\parindent}{0pt}
\setlength{\parskip}{0.47em}
\emergencystretch=2em
\lstset{basicstyle=\ttfamily\footnotesize,breaklines=true,frame=single,columns=fullflexible,showstringspaces=false}
\pagestyle{fancy}\fancyhf{}
\lhead{Game Design \& Engineering}\rhead{Working Example and Guidance}\cfoot{\thepage}
\begin{document}
{\huge\bfseries\color{suiteblue} Game Balance Experiment -- Template}\par
\vspace{0.15in}
{\normalsize\color{suitegray} Illustrative examples are fictional and must be replaced with actual project evidence}\par
\vspace{0.15in}\hrule\vspace{0.15in}

\section{Instructions}
Use the worked example to understand the required level of specificity. Keep identifiers stable across related design records. Replace each example with your project decision and justify any omitted field. Examples are fictional and not findings reported by Schell.


\section{Document control}
\begin{tabularx}{\linewidth}{@{}lX@{}}\toprule Record version & [Version identifier]\\ Owner & [Accountable owner]\\ Review date & [Review date]\\ Project & [Project name]\\\bottomrule\end{tabularx}

\subsection{Design question}
Example: Does adding a low-damage high-accuracy tool create a meaningful choice?

Your project: [Enter a specific value, justification, owner, and evidence link.]


\subsection{Competing strategies}
Example: Precise tool versus high-damage unreliable tool

Your project: [Enter a specific value, justification, owner, and evidence link.]


\subsection{Outcome measure}
Example: Expected damage and success rate across three enemy profiles

Your project: [Enter a specific value, justification, owner, and evidence link.]


\subsection{Parameter range}
Example: Accuracy from 0.50 to 0.95; damage from two to six

Your project: [Enter a specific value, justification, owner, and evidence link.]


\subsection{Control condition}
Example: Baseline enemy at health ten with no armor

Your project: [Enter a specific value, justification, owner, and evidence link.]


\subsection{Experiment seeds}
Example: Record five named seeds and the PRNG implementation

Your project: [Enter a specific value, justification, owner, and evidence link.]


\subsection{Sample size}
Example: At least ten thousand simulated encounters per condition

Your project: [Enter a specific value, justification, owner, and evidence link.]


\subsection{Interpretation limits}
Example: Expected damage does not measure perceived fun or player skill

Your project: [Enter a specific value, justification, owner, and evidence link.]


\subsection{Observed result}
Example: Illustrative only: precision tool has lower variance

Your project: [Enter a specific value, justification, owner, and evidence link.]


\subsection{Decision}
Example: Retain both tools pending player comprehension testing

Your project: [Enter a specific value, justification, owner, and evidence link.]


\section{Approval and traceability}
Requirement ID: [Identifier]. Reviewer: [Reviewer]. Status: [Draft/Reviewed/Approved]. Supporting evidence: [Evidence identifier].


\end{document}
